\section{Procedencia de los resultados y reproducción de la ampliación}
\label{ap:procedencia-reproduccion}
\begingroup
\setlength{\parskip}{1pt}

La exposición parte del núcleo APP--TRIT--TPK y mantiene su construcción
correlativa del continuo. Las lecturas regionales, el producto nativo y los
relojes irreducibles preceden a sus realizaciones analíticas. La integración
de centro, frontera y memoria reúne resultados del corpus HMT--MD con
las composiciones desarrolladas para el doble círculo. La procedencia se
registra por resultado, de modo que una formulación reunida o un certificado
adicional no convierte retrospectivamente en nueva la estructura que utiliza.

\subsection{Dependencias reunidas en el cuerpo}

El núcleo comprende las definiciones y operaciones comunes, junto con
sus dos figuras. Las definiciones, enunciados y pruebas de los cinco
desarrollos anteriores se conservan íntegramente.

La coordenatización central, su frontera levantada y las firmas de transporte se
desarrollan en la sección~\ref{sec:centro-frontera}. La distinción entre
el punto central \((5,5)\), la firma producto \(555555\) y las hojas que
su transporte distingue se conserva en la definición~\ref{def:firma555}
y la proposición~\ref{prop:firma555}. El calendario de capacidad y sus
dos niveles de retorno se prueban en los
teoremas~\ref{vac-add-teorema-primario} y~\ref{vac-add-teorema-secundario}.
La geometría de la coordenatización no se identifica con la historia completa: su
recuperación utiliza la ley multiafín y los cocientes conservados en la frontera.

El transporte del doble círculo y los momentos de memoria se reúnen en
la sección~\ref{sec:transporte-momentos}. Los lectores de cola gamma de
la sección~\ref{sec:lectores-gamma} reconstruyen las columnas con sus
dominios y productos cruzados. La sección~\ref{sec:schur-coercividad}
establece la estimación en una ventana completa y su transporte a
subespacios con colas. Las igualdades de conservación y la positividad
de una forma son afirmaciones distintas: cada resultado de esta ampliación
conserva el dominio en el que se demuestra.

Los antecedentes formales necesarios para estas composiciones se
desarrollan en los apéndices A--C: reglas de emisión e historias compatibles,
lecturas regionales y retorno nonádico, y registro dodecafásico reversible.
Las ecuaciones y pruebas utilizadas permanecen así expuestas junto con
los dominios y datos de representación que requiere cada resultado.

\clearpage
\subsection{Orden de reproducción y alcance de los controles}

La edición conserva los antecedentes, los cinco desarrollos anteriores y
los controles focales que acompañan sus identidades. Su compilación acredita
una realización material legible del texto; la conservación documental y
las comprobaciones computacionales no certifican por extensión todos los
resultados de un proyecto. La aportación de la ampliación consiste en
incorporar las composiciones siguientes con sus demostraciones.

La carpeta que acompaña al PDF conserva las fuentes \LaTeX{}, los cinco
desarrollos en Markdown, los programas y los recibos JSON. La reproducción
sigue la dependencia matemática de los objetos:

\begin{enumerate}
\item Construir los lectores de centro y frontera conservando residuo y
cociente, y comprobar la partición de las 324 semillas bajo los transportes
de avance y media vuelta.
\item Evaluar las partes enteras del calendario de vacancias mediante
horquillas racionales. Las cotas con resto preceden a las cifras y fijan
la certificación de cada posición.
\item Comprobar los balances de costura, los terminales de memoria y las
factorizaciones de momentos; conservar los términos cruzados al formar
las matrices conjuntas.
\item Evaluar las cotas locales de cola gamma y del complemento de Schur,
con los dominios de las estimaciones impresos en sus enunciados.
\item Compilar el manuscrito y cotejar las fuentes comunes, las ecuaciones
trasladadas, las referencias y la disposición de las páginas.
\end{enumerate}

Los controles racionales y los intervalares verifican las identidades y
las cotas finitas que declaran. Las demostraciones funcionales establecen
las afirmaciones universales correspondientes. La conservación de los
archivos y una compilación correcta son controles editoriales separados;
ninguno de ellos se presenta como sustituto de una demostración global.

La ejecución ordinaria y la optimizada de los programas utiliza condiciones
explícitas: desactivar las aserciones del intérprete no altera el criterio
de aceptación. El archivo de instrucciones de la entrega especifica los
comandos y las dependencias del entorno. Los originales anteriores se
conservan por separado y no son sobrescritos por esta ampliación.
\par\endgroup
